\chapter{Renin}

\section*{What Is This Test?}
Renin is an aspartyl protease enzyme produced, stored, and secreted by the juxtaglomerular (JG) cells of the afferent arterioles of the kidney. Renin is synthesized as a preprorenin precursor that is cleaved to prorenin and then to active renin, which is stored in JG cell granules. Renin secretion is regulated by three primary mechanisms: (1) renal baroreceptor sensing of decreased afferent arteriolar perfusion pressure, (2) sodium chloride delivery to the macula densa cells of the distal convoluted tubule (decreased NaCl delivery stimulates renin), and (3) sympathetic nervous system activation via beta-1 adrenergic receptors on JG cells. Renin cleaves angiotensinogen (produced by the liver) into angiotensin I, which is then converted to angiotensin II by angiotensin-converting enzyme (ACE) primarily in the pulmonary endothelium. Angiotensin II is a potent vasoconstrictor and stimulates aldosterone secretion from the adrenal zona glomerulosa. The renin-angiotensin-aldosterone system (RAAS) thus regulates blood pressure, sodium and potassium balance, and extracellular fluid volume. The renin test is measured as either plasma renin activity (PRA), which measures the enzymatic generation of angiotensin I from endogenous angiotensinogen per unit time (ng/mL/hr), or as direct renin concentration (DRC), which measures immunoreactive active renin mass by immunoassay (mU/L or pg/mL) \cite{tietz2012textbook}. The renin test is essential for interpreting aldosterone levels, calculating the aldosterone-to-renin ratio (ARR) for primary aldosteronism screening, and evaluating hypertensive disorders.

\section*{Alternative Names}
\begin{itemize}
  \item Plasma Renin Activity (PRA)
  \item Direct Renin Concentration (DRC)
  \item Active Renin
  \item Serum Renin
  \item Renin Activity, Plasma
  \item Aldosterone-to-Renin Ratio (ARR)
  \item Renal Vein Renin Sampling
\end{itemize}
\section*{Principle}
Plasma renin activity (PRA) is measured by a two-step process: (1) the patient's plasma (containing renin and angiotensinogen) is incubated at 37 degrees C for a defined time (typically 90 minutes to 3 hours) at pH 5.5--7.4, during which renin enzymatically generates angiotensin I from endogenous angiotensinogen, in the presence of angiotensin-converting enzyme and angiotensinase inhibitors (EDTA, PMSF, 8-hydroxyquinoline) to prevent degradation of generated angiotensin I; (2) the angiotensin I generated is measured by competitive radioimmunoassay (RIA) or enzyme immunoassay (EIA). PRA is reported as nanograms of angiotensin I generated per milliliter per hour (ng/mL/hr or ng/mL/h). PRA is a functional (enzymatic) assay, not an immunoassay for renin mass. Direct renin concentration (DRC) is measured by a two-site sandwich immunometric (chemiluminescent) immunoassay using monoclonal antibodies against active renin. DRC reports the immunoreactive mass of active renin in mU/L or pg/mL (1 pg/mL active renin = approximately 1.6 microunits/mL). DRC is performed on automated platforms (Roche Cobas, Diasorin Liaison, Siemens IMMULITE) and is simpler and faster than PRA but less extensively validated in published literature for primary aldosteronism screening. Conversion between PRA and DRC is possible but not linear (PRA depends on endogenous angiotensinogen levels, which vary with estrogen, pregnancy, liver disease, nephrotic syndrome). PRA is the gold standard for ARR screening per Endocrine Society guidelines, though DRC is increasingly used because of its automation and reproducibility.

\section*{History}
Renin was discovered in 1898 by Robert Tigerstedt and Per Bergman at the Karolinska Institute in Stockholm, who demonstrated that injection of rabbit kidney extract into rabbits raised blood pressure and named the substance "renin." \cite{tigerstedt1898niere} The RAAS was further elucidated by Irvine Page, Eduardo Braun-Menéndez, and others in the 1930s--1940s. Page and Helmer discovered the renin substrate (angiotensinogen) in 1940. Angiotensin was isolated and characterized in the 1950s by Skeggs and colleagues. The radioimmunoassay for angiotensin I and measurement of plasma renin activity was developed by Haber and colleagues in 1969 \cite{haber1969application}. The aldosterone-to-renin ratio (ARR) as a screening test for primary aldosteronism was first described by Hiramatsu in 1981 and subsequently validated by many groups \cite{young2007primary}. Direct renin concentration (DRC) immunoassays were developed in the 1990s. The Endocrine Society clinical practice guidelines by Funder and colleagues (2008, updated 2016) established the ARR as the screening test of choice for primary aldosteronism \cite{funder2016management}. The development of renin inhibitors (aliskiren) as a therapeutic drug class in the 2000s provided a pharmacologic proof of concept for the RAAS in hypertension.

\section*{Why This Test Is Ordered}
\begin{itemize}
  \item Calculation of the aldosterone-to-renin ratio (ARR) for screening for primary aldosteronism in hypertensive patients --- the single most common use of renin testing
  \item Classification of hypertension as low-renin (primary aldosteronism, Liddle syndrome, apparent mineralocorticoid excess, licorice ingestion) vs. high-renin (renovascular hypertension, renin-secreting tumors, malignant hypertension)
  \item Identification of secondary hyperaldosteronism in hypertension with hypokalemia --- renal artery stenosis (atherosclerotic or fibromuscular dysplasia), renin-secreting tumors (reninomas)
  \item Investigation of unexplained hypokalemia with metabolic alkalosis and normal or low blood pressure --- Bartter syndrome and Gitelman syndrome (salt-wasting nephropathies with secondary hyperreninemic hyperaldosteronism)
  \item Classification of adrenal insufficiency (primary vs. secondary) --- elevated renin due to aldosterone deficiency in primary adrenal insufficiency
  \item Renal vein renin sampling in suspected renovascular hypertension --- lateralized renin secretion (ratio >1.5:1 affected vs. contralateral kidney) indicates renin-dependent hypertension potentially curable by revascularization
  \item Research in cardiovascular risk stratification --- elevated renin is independently associated with increased risk of myocardial infarction, stroke, and mortality in hypertensive populations
  \item Monitoring medical therapy for primary aldosteronism --- renin should rise into the normal range during adequate mineralocorticoid receptor antagonist dosing (spironolactone/eplerenone), indicating blockade of the MR and restoration of normal RAAS feedback
\end{itemize}

\section*{Primary vs Secondary Test}
Renin measurement alone is rarely of diagnostic value; it is almost always interpreted in conjunction with aldosterone as the aldosterone-to-renin ratio (ARR), which is a primary screening test for primary aldosteronism. Isolated renin measurement is most useful for classification of abnormal potassium and blood pressure states. The ARR is recommended by the Endocrine Society as the screening test for primary aldosteronism in at-risk hypertensive populations \cite{funder2016management}.

\section*{How This Test Is Performed}
For PRA, a blood sample is drawn from a vein into a pre-chilled EDTA (purple-top) tube. The tube must be inverted gently (not shaken) to mix anticoagulant. The specimen is kept at room temperature (NOT on ice, because cryoactivation of prorenin to renin at 0--4 degrees C can falsely elevate PRA) and centrifuged at room temperature within 30 minutes. Plasma is separated and frozen at -20 degrees C or colder. For DRC, blood is collected into an EDTA tube and plasma is separated within 4 hours and can be stored at 2--8 degrees C for 24 hours or frozen for longer. The patient should be in a controlled posture at the time of blood draw. For a standardized ARR: blood is drawn mid-morning (8--10 AM) after the patient has been upright (sitting, standing, walking) for at least 2 hours, then seated for 5--15 minutes before venipuncture. Specimens collected after prolonged supine rest will have lower renin values. For renal vein renin sampling: catheters are placed in both renal veins via femoral vein access; renin is measured from both renal veins and the inferior vena cava below the renal veins; a renal vein renin ratio (affected:unaffected) >1.5:1 suggests renin-dependent hypertension that may benefit from revascularization.

\section*{What Preparation Is Needed}
Patient preparation is extensive and critical for ARR interpretation: the patient should be on a liberal sodium diet (not sodium restricted) for at least 3 days before testing because sodium restriction stimulates RAAS and can cause a false-positive ARR. Hypokalemia must be corrected to >4.0 mmol/L with oral potassium supplementation because hypokalemia suppresses aldosterone, causing false-negative ARR. Several medications must be discontinued or substituted before ARR: mineralocorticoid receptor antagonists (spironolactone, eplerenone) --- hold for 4--6 weeks, as these directly and persistently elevate renin; high-dose amiloride and triamterene --- hold for 2--4 weeks; loop and thiazide diuretics --- hold for 4 weeks (stimulate renin, can cause false-negative ARR); ACE inhibitors, ARBs, direct renin inhibitors (aliskiren) --- hold for 2 weeks if possible (decrease aldosterone, increase renin, can cause false-negative ARR); beta-blockers, clonidine, alpha-methyldopa --- hold for 2 weeks if possible (suppress renin, can cause false-positive ARR); NSAIDs --- hold for 2 weeks if possible (suppress renin); oral contraceptives and estrogen therapy affect angiotensinogen and PRA but not DRC --- document their use. Patients who cannot safely discontinue antihypertensives can be screened while on non-interfering agents: verapamil slow-release, hydralazine, prazosin, or doxazosin. Fasting for 8--10 hours is recommended. The patient must be upright for at least 2 hours before venipuncture and then seated for 5--15 minutes, with position and time documented because renin is profoundly affected by posture.

\section*{Contraindications and Precautions}
\begin{itemize}
  \item Pre-analytical factors are major sources of error: PRA must NOT be collected on ice because cold activates prorenin to active renin (cryoactivation), artificially doubling or tripling measured renin activity. Specimens for PRA should be kept at room temperature, centrifuged at room temperature within 30 minutes, and plasma frozen immediately. For DRC, specimens should be maintained at room temperature; cryoactivation affects DRC to a lesser degree but still occurs. Hemolyzed specimens are unacceptable because hemoglobin interferes with angiotensin I generation and immunoassays. Prolonged tourniquet use and fist clenching during venipuncture decrease blood flow and may alter renin release. Strenuous exercise 24 hours before testing should be avoided (sympathetic activation elevates renin). Renin levels fall progressively with age, particularly after age 65, and ARR increases --- older patients have a higher false-positive rate. Chronic kidney disease (CKD stages 3--5) is associated with low renin and high ARR, making ARR interpretation less reliable. African ancestry is associated with lower renin levels compared to Caucasians. Factors that decrease PRA/DRC: high-sodium diet, volume expansion, mineralocorticoid excess, beta-blockers, NSAIDs, clonidine, advanced age, CKD, African ancestry, autonomous aldosterone secretion. Factors that increase PRA/DRC: low-sodium diet, volume depletion, diuretics, upright posture, pregnancy (estrogen stimulates angiotensinogen, raising PRA but not necessarily DRC), ACE inhibitors/ARBs, renovascular hypertension, Bartter syndrome/Gitelman syndrome, cirrhosis with ascites, congestive heart failure. PRA is specifically elevated by high angiotensinogen states (pregnancy, estrogen therapy, oral contraceptives), while DRC is less affected because it measures renin mass, not activity.
\end{itemize}
\section*{Risks}
Minimal risk from venipuncture: slight pain, bruising, or hematoma. From renal vein sampling: groin hematoma (common), contrast reaction, renal vein dissection, and rare renal infarction. Renal vein sampling is a specialized procedure performed at experienced centers only.

\section*{What Happens After the Test}
Results are available within 1--3 days (PRA is often a send-out test; DRC may be available sooner). The ARR is calculated: if the ARR is elevated (PRA-based ARR >30, or DRC-based ARR >3.7) AND aldosterone is >10--15 ng/dL, the screen is positive. A positive screen must be confirmed by repeat ARR and then by a confirmatory test (saline suppression test, captopril challenge, oral sodium loading, or fludrocortisone suppression) before the diagnosis of primary aldosteronism is established. If primary aldosteronism is confirmed, subtype differentiation (unilateral adenoma vs. bilateral hyperplasia) requires adrenal CT and, for surgical candidates, adrenal vein sampling. If renin is elevated with elevated aldosterone and a normal ARR, secondary hyperaldosteronism is present. Renal artery Doppler ultrasound or CT/MR angiography should be considered for renovascular hypertension. If renin is low with low aldosterone, other causes of hyporeninemic hypoaldosteronism should be investigated (including diabetic nephropathy and other forms of Type IV renal tubular acidosis). During medical therapy for primary aldosteronism with mineralocorticoid receptor antagonists, renin should rise into the normal range once the MR is adequately blocked; persistent suppression despite treatment indicates inadequate dosing.

\section*{Understanding Results}

\subsection*{Below Normal}
\begin{itemize}
  \item \textbf{Suppressed renin (PRA <1.0 ng/mL/hr or DRC <10 mU/L) with elevated aldosterone:} Primary aldosteronism --- most common cause of suppressed renin. Autonomous aldosterone secretion from an adrenal adenoma or bilateral hyperplasia expands extracellular volume, suppresses renin via the renal baroreceptor mechanism, and produces hypertension with variable hypokalemia.
  \item \textbf{Low renin with normal or low aldosterone:} Low-renin (volume-dependent) essential hypertension --- the most common form of essential hypertension, accounting for 25--30\% of hypertensive patients. Normal ARR because both renin and aldosterone are proportionately low. Associated with sodium sensitivity, older age, African ancestry, and good response to diuretics.
  \item \textbf{Low renin with low aldosterone and hyperkalemia:} Hyporeninemic hypoaldosteronism (Type IV RTA) --- low renin due to juxtaglomerular apparatus damage in diabetic nephropathy, chronic interstitial nephritis, NSAID use, or in older adults with mild renal insufficiency. Hyperkalemia and mild metabolic acidosis (normal anion gap) are characteristic.
  \item \textbf{Low renin with low aldosterone and hypertension:} Liddle syndrome (gain-of-function mutation in ENaC, autosomal dominant) --- constitutive sodium reabsorption in collecting duct drives volume expansion, hypertension, hypokalemia, and metabolic alkalosis. Renin and aldosterone are both suppressed. Differentiated from primary aldosteronism by low ARR (aldosterone also suppressed) and response to amiloride/triamterene but not to spironolactone. Apparent mineralocorticoid excess (11-beta-HSD2 deficiency, genetic or licorice ingestion) --- cortisol cannot be converted to cortisone, allowing cortisol to activate the mineralocorticoid receptor. Renin and aldosterone are suppressed, hypertension and hypokalemia occur.
  \item \textbf{Suppressed renin from licorice/glycyrrhizic acid ingestion:} Glycyrrhetinic acid in licorice inhibits 11-beta-HSD2, causing apparent mineralocorticoid excess. Suppressed renin and aldosterone with hypertension and hypokalemia. Reversible with cessation.
  \item \textbf{Patients on beta-blockers:} Beta-1 adrenergic receptor blockade suppresses renin secretion from JG cells. ARR may be falsely elevated.
\end{itemize}

\subsection*{Above Normal}
\begin{itemize}
  \item \textbf{Elevated renin with elevated aldosterone (secondary hyperaldosteronism):} Renovascular hypertension due to atherosclerotic renal artery stenosis (>70\% stenosis, typically ostial lesion in older adults with cardiovascular risk factors) or fibromuscular dysplasia (younger women, beaded appearance of mid-distal renal artery on angiography). Renal hypoperfusion activates JG cell renin release, driving RAAS, causing hypertension that is renin-dependent and potentially curable by revascularization (stenting for atherosclerosis, angioplasty for FMD). Renal vein renin ratio >1.5:1 predicts blood pressure improvement after revascularization \cite{rimoldi2014secondary}.
  \item \textbf{Renin-secreting juxtaglomerular cell tumor (reninoma):} Rare benign renal tumor secreting excess renin. Young adults (typically 20--30 years old). Severe hypertension, hypokalemia, and markedly elevated renin and aldosterone. CT/MRI shows small (<2--3 cm) renal cortical tumor. Surgical resection cures hypertension.
  \item \textbf{Bartter syndrome and Gitelman syndrome:} Inherited renal tubular salt-wasting disorders causing chronic volume depletion. Bartter syndrome (defects in thick ascending limb transporters NKCC2, ROMK, CLC-Kb) --- neonatal/severe presentation, hypercalciuria, normal or low serum magnesium. Gitelman syndrome (defect in distal convoluted tubule NCC transporter) --- later presentation, hypocalciuria, hypomagnesemia. Both cause hypokalemic metabolic alkalosis, elevated renin and aldosterone, and normal to low blood pressure (NOT hypertensive, distinguishing from renovascular disease and primary aldosteronism).
  \item \textbf{Volume depletion states:} Diuretic therapy (thiazides, loop diuretics), vomiting, diarrhea, excessive sweating, low-sodium diet --- renin elevated by renal baroreceptor, aldosterone appropriately elevated. Normal ARR.
  \item \textbf{Congestive heart failure and cirrhosis with ascites:} Reduced effective circulating volume (despite total body sodium and water excess) activates RAAS. Elevated renin and aldosterone contribute to sodium and water retention. Prognostic marker --- markedly elevated renin/aldosterone levels predict worse outcomes in heart failure.
  \item \textbf{Malignant hypertension (accelerated phase):} Markedly elevated blood pressure (often >200/130 mmHg) with renal arteriolar damage and ischemia, driving massive renin release. Fundoscopy shows grade III or IV hypertensive retinopathy. Renal function often impaired. Renin and aldosterone both markedly elevated.
  \item \textbf{Pregnancy:} Estrogen stimulates hepatic angiotensinogen production, increasing PRA disproportionately (PRA increased 3--4 fold). DRC is less affected. Physiologic RAAS activation during pregnancy maintains blood pressure despite progesterone-induced vasodilation.
\end{itemize}

\section*{Normal Results}
\begin{itemize}
  \item \textbf{Plasma renin activity (PRA), supine:} 0.2--1.6 ng/mL/hr; upright (2 hr ambulation): 1.5--5.2 ng/mL/hr --- values 2--4 times higher upright than supine
  \item \textbf{PRA, sodium-restricted diet:} Upright PRA can increase to 5--10 ng/mL/hr
  \item \textbf{Direct renin concentration (DRC), supine:} 3--16 mU/L (males), 2--16 mU/L (females); upright: 10--40 mU/L
  \item \textbf{PRA in primary aldosteronism:} <1.0 ng/mL/hr (often <0.5 ng/mL/hr, undetectable)
  \item \textbf{PRA in renovascular hypertension:} >5 ng/mL/hr (often markedly elevated >10--20 ng/mL/hr)
  \item \textbf{Aldosterone-to-renin ratio (ARR) using PRA:} <20 ng/dL per ng/mL/hr (normal); 20--30 (indeterminate); >30 (positive screen for primary aldosteronism)
  \item \textbf{ARR using DRC:} <2 (pmol/L per mU/L) normal; >3.7 positive screen
  \item \textbf{Renal vein renin ratio (affected:unaffected kidney):} >1.5:1 suggests lateralized renin secretion; >2:1 very strongly lateralized
\end{itemize}

\section*{Follow-up Tests}
\begin{itemize}
  \item \textbf{Simultaneous aldosterone and ARR calculation:} Renin must always be interpreted with aldosterone. The ARR is the first step in primary aldosteronism screening.
  \item \textbf{Saline suppression test:} After positive ARR, 2 L normal saline IV over 4 hours. Failure of aldosterone to suppress <5--10 ng/dL confirms primary aldosteronism.
  \item \textbf{Captopril challenge test:} Captopril 25--50 mg orally. Aldosterone >12--15 ng/dL and ARR remaining elevated confirms primary aldosteronism.
  \item \textbf{Oral sodium loading test:} 5000 mg sodium diet for 3 days with 24-hour urine aldosterone >12--14 mcg/24 hr on high sodium diet confirms primary aldosteronism.
  \item \textbf{Adrenal CT (fine-cut, 1--3 mm):} Detects adrenal adenoma vs. bilateral hyperplasia. Determines presence of mass but NOT functionality.
  \item \textbf{Adrenal vein sampling (AVS):} Gold standard for lateralization. Aldosterone/cortisol ratio dominant side >4:1 with contralateral suppression confirms unilateral adenoma \cite{rossi2014expert}.
  \item \textbf{Renal artery Doppler ultrasound, CT angiography, or MR angiography:} For suspected renal artery stenosis. Duplex ultrasound with peak systolic velocity >180--200 cm/s or renal-to-aortic ratio >3.5 suggests >60\% stenosis \cite{rimoldi2014secondary}.
  \item \textbf{Renal vein renin sampling:} Considered when noninvasive imaging is equivocal for renal artery stenosis functional significance. Lateralized renin ratio >1.5:1 predicts BP improvement after revascularization.
  \item \textbf{Genetic testing:} For familial primary aldosteronism (CYP11B1/CYP11B2 chimeric gene for GRA Type I, KCNJ5 mutations for Type III). For Liddle syndrome (SCNN1B, SCNN1G mutations). For apparent mineralocorticoid excess (HSD11B2 mutations).
  \item \textbf{Urine potassium (24-hour):} Inappropriately high urinary potassium (>30 mmol/24 hr) in the setting of hypokalemia confirms renal potassium wasting, characteristic of primary and secondary hyperaldosteronism.
\end{itemize}

\section*{References}
\bibliographystyle{unsrtnat}
\bibliography{references}

\clearpage
\section*{Regulatory Status and Authorization Date}
This chapter describes a test category, not a specific commercial product. FDA, CDSCO, EU IVDR, and MHRA approval or registration dates therefore cannot be assigned without the assay manufacturer, product name, instrument, and intended use. For laboratory-developed tests, an individual agency approval date may not exist. See the Preface for the interpretation of regulatory status.

\clearpage
